% =====================================================================
% Section 3 — Microcode as an Implementational Layer
% Draft for TCHES submission.
%
% Scope: this section introduces the microcode layer as a programming
% target and the GENERAL methodology by which we compile a cryptographic
% primitive to a patch that runs on a stock Intel Goldmont part. It first
% explains, for a reader who has not met them, exactly what a triad is and
% what the patch RAM is, building on the reverse engineering of Koppe et
% al. (AMD K8/K10) and Borrello et al. (CustomProcessingUnit, Intel
% Goldmont) and on Krog's lib-micro (https://libmicro.dev), the C macro
% language / patching library whose infrastructure this project actually
% uses (cf. EXPERIMENTS.md, include/patch.h: patch_ucode(),
% hook_match_and_patch()). The triad/sequence-word bit-field encoding in
% sec:triad-anatomy follows libmicro.dev/structure.html. It then describes
% how that mechanism is used to customise an
% instruction, how triads are packed, how the 128-triad budget shapes a
% design, the toolchain that lowers an op-list to a verified patch, and
% the undocumented hardware behaviour we had to map. No primitive is
% implemented here. The two worked examples follow: Keccak-f[1600] in
% Section 4 (control flow and register residency) and the Curve25519
% field arithmetic in Section 5 (dense multiply-accumulate-carry).
%
% Cross-reference contract with later sections: Sections 4 and 5 refer to
% the labels sec:ucode-layer, sec:triads-packing, sec:budget and
% sec:findings; these are preserved and their content still matches the
% citing text (e.g. sec:triads-packing carries the two-macroinstruction
% hook and the generator discipline).
% =====================================================================

\section{Microcode as an Implementational Layer}
\label{sec:ucode-layer}

Prior work has studied microcode as an object to be reverse engineered and,
latterly, as a surface on which to mount security mechanisms. We instead treat
it as a programming target: a layer beneath the instruction set on which a
cryptographic primitive can be expressed directly, and on which it can be made
to run faster than any code the instruction set itself admits. This section
introduces the layer in the terms a reader will need, namely what a triad is,
what the writable patch RAM is, and by what mechanism a stock instruction is
customised to run code of our choosing, and it then sets out the methodology
that turns those ingredients into a verified patch on the Intel
Goldmont~\hbox{Celeron~N3350}. The same methodology underlies every primitive
in this work, namely the field arithmetic for P-256, Curve25519, secp256k1 and
Poly1305 together with the Keccak-$f[1600]$ permutation of \Cref{sec:keccak},
and we therefore present it once, in general terms, and instantiate it twice in
the sections that follow.

Two properties of the layer motivate the entire effort, and both will recur as
the source of every result we report. The first is that the
micro-architectural integer register file is twice the size of the
architectural one, comprising 16 general-purpose and 16 temporary registers,
which lets a primitive keep far more state resident than the roughly 14 usable
general-purpose registers of native x86 and thereby removes the register spills
that dominate compiled code on a register-starved in-order core; this is the
property that \Cref{sec:keccak} exploits to hold an entire Keccak state in
registers across all 24 rounds. The second is that a triad issues up to three
micro-ops as a single sequential bundle, which yields a denser and more
hazard-aware packing than a linear instruction stream affords, and in
particular lets several carry chains advance at once; this is the property that
\Cref{sec:curve25519} exploits to overlap the multiply--accumulate--carry chain
of a reduced field multiplication. The price of admission is a hard budget of
128 triads, together with a body of hardware behaviour that is nowhere
documented and that we were obliged to map empirically.

% ---------------------------------------------------------------------
\subsection{Microcode, Triads, and the Patch RAM}
\label{sec:triad-anatomy}

\paragraph{From macroinstructions to micro-ops.}
The x86 instruction set is a contract between software and the hardware that
runs it, a stable interface that fixes what each instruction must appear to do
while saying nothing about how the hardware is to achieve it, and it is not the
language the core actually executes. A decode unit translates each architectural instruction, which we
call a \emph{macroinstruction}, into one or more internal \emph{micro-ops}
($\mu$ops) drawn from a simpler, RISC-like repertoire that the execution units
actually run. Frequent and simple macroinstructions are handled by fast
hardware decoders that emit a handful of $\mu$ops directly, whereas complex or
rare ones are instead decoded by a microcode sequencer. It is for these latter,
microcoded instructions that the $\mu$op program implementing the
macroinstruction is held on an on-chip microcode ROM, the MSROM, from which the
sequencer fetches it. Goldmont exposes more
than two thousand distinct $\mu$ops, and a single complex macroinstruction may
expand to dozens of them. The encoding is \emph{vertical}, meaning that each
$\mu$op resembles a compact RISC instruction with an operation field and a
small number of operand fields rather than a wide horizontal control word. The
reverse engineering of Koppe et al.\ \cite{koppeReverseEngineeringX862017} on
AMD K8/K10 first established this structure publicly, organising the $\mu$ops
into operation classes for register arithmetic, loads, stores and special
operations, and Borrello et al.\ \cite{borrelloCustomProcessingUnitReverseEngineering2023}
subsequently recovered the corresponding $\mu$ops for Goldmont together with
their operand and immediate encodings.

\paragraph{The triad.}
Inside the microcode store the $\mu$ops are not laid out singly but grouped into
\emph{triads}. A triad is a bundle of three $\mu$ops together with one
\emph{sequence word}, and it is the unit in which microcode is addressed,
fetched, and issued: the sequencer advances triad by triad, presenting the
three $\mu$ops of a triad to the execution units before it moves on. The
sequence word is the control-flow field of the triad. It selects what the
sequencer does after the triad's $\mu$ops have issued, namely fall through to
the next triad, branch to a named triad address, or signal that decoding of the
macroinstruction is complete, and on Goldmont it can additionally impose a
synchronisation barrier on the bundle, in the manner of an \texttt{lfence}.
A triad is thus at once a packing unit and the granularity of control flow. It
is a packing unit because its three $\mu$ops are fetched and issued together as
a single bundle, so that the quantity governing both the size of a program and
the length of its dependency chains is the number of triads rather than the
number of $\mu$ops, and filling all three slots of every triad is in consequence
the central concern of the packing described in \Cref{sec:triads-packing}. It is
the granularity of control flow because the sequence word steers the sequencer
one triad at a time and a branch can target only a triad boundary, never a
$\mu$op within a triad, which has the practical consequence, developed in
\Cref{sec:findings}, that an operation setting a condition and the branch that
tests it must be placed in the same triad. The MSROM on Goldmont holds on the
order of eight thousand such triads, 7936 to be exact
\cite{borrelloCustomProcessingUnitReverseEngineering2023}, which together form
the $\mu$op program that decodes the entire x86 instruction set. Against that backdrop the 128
triads of writable patch RAM introduced below amount to under two percent of the
decode store, and it is within that small allowance that an entire cryptographic
primitive must be expressed.

\paragraph{The encoding of a triad.}
The structure of a triad is worth making concrete, since it determines what a
single $\mu$op can express and therefore what the packer of \Cref{sec:findings}
is working against. We follow the Goldmont encoding first documented by Ermolov,
Goryachy and Sklyarov's uCodeDisasm \cite{ermolovUCodeDisasm2021}, whose $\mu$op
mnemonics our toolchain inherits by way of lib-micro. Each $\mu$op is a 48-bit
word. A 12-bit opcode names the operation, and for some operations also encodes
a sub-mode such as a condition code or a data size; two 6-bit fields select the
source registers and a third 6-bit field the destination, which for a memory
operation instead supplies a second source. An operand need not be a register:
when the top three bits of a source field are \texttt{001} that field's
remaining bits and the $\mu$op's two immediate fields together form a 16-bit
immediate, and when they are \texttt{010} the field indexes a read-only
constants ROM. A few mode bits and a two-bit parity check complete the word. The
16-bit ceiling on an immediate is exactly why a 64-bit constant must be built by
composition, a point \Cref{sec:findings} returns to. The sequence word is a
narrower 30-bit word, but its fields act on the triad rather than on data.
Three 2-bit pointers bind its actions to particular slots; a four-bit
control-flow field selects among returning to the caller, looping within the
microcode, stopping, and signalling that decoding is complete; a branch-address
field supplies a jump target that becomes conditional when the preceding $\mu$op
is a flag-setting test or subtract; and a three-bit synchronisation field can
place a load fence or a full barrier on the bundle, which is the
\texttt{lfence}-like behaviour noted above. These two fields together are the
whole of control flow in the layer, and \Cref{sec:findings} records which of
their forms behave reliably on this part.

\paragraph{The patch RAM.}
The MSROM is fixed at manufacture, but a CPU must be able to correct decode
bugs after it ships. Vendors provide for this with a small writable microcode
store, the MSRAM or \emph{patch RAM}: in normal operation the firmware or
operating system loads a cryptographically signed microcode patch into it at
boot by writing the patch's address to a dedicated model-specific register, and
the store is volatile and must be reloaded after every reset. It is one of
several microcode-sequencer arrays, the others holding the read-only $\mu$ops,
their sequence words, and the match table described below. On Goldmont the patch
RAM occupies the region \texttt{U7c00}--\texttt{U7e00} and holds at most 128
triads, and this ceiling, as \Cref{sec:budget} develops at length, is the
binding resource of our entire approach. A patch is therefore a short array of
triads, each consisting of three $\mu$ops and a sequence word, written in
exactly the format the MSROM itself uses.

\paragraph{Customising an instruction.}
A patch is inert until some macroinstruction is made to decode into it. This is
the role of the \emph{match-and-patch} mechanism, a table of 32 entries each of
which redirects a read-only triad address to one in the patch RAM. When the
sequencer is about to fetch the triad at a matched MSROM address it is instead
sent to the corresponding triad in the MSRAM. Pointing a match entry at the
MSROM entry point of a chosen macroinstruction therefore \emph{hooks} that
instruction, so that issuing it transfers control into our patch. Because each
entry redirects an even-aligned pair of addresses at once, mapping both
\texttt{src} and \texttt{src}${+}1$ to consecutive patch triads, only even
addresses can be hooked. This is the mechanism, and the only mechanism, by which
we install code of our own: we write triads into the patch RAM and aim a match
entry at an instruction we are willing to repurpose.

Three lines of prior work make this practical on Goldmont, and our toolchain
rests on all of them. The foundation is due to Ermolov, Goryachy and Sklyarov:
beginning in 2020 they obtained Intel's \emph{Red Unlock} debug state on Goldmont
through the Intel-SA-00086 Management Engine vulnerability, identified the
undocumented \texttt{udbgrd} and \texttt{udbgwr} instructions that read and write
the microcode arrays from software \cite{ermolovUndocumentedX862022}, and released
\emph{uCodeDisasm}, the disassembler that first documented the Goldmont $\mu$op,
triad and sequence-word encoding publicly \cite{ermolovUCodeDisasm2021}. Building
on this, Borrello et al.'s CustomProcessingUnit (2023) reverse-engineered the
patch and hook format, defined the semantics of some eight thousand $\mu$ops, and
provides a static decompiler together with an assembler from $\mu$op mnemonics to
triads \cite{borrelloCustomProcessingUnitReverseEngineering2023}. We build most
directly on Krog's \emph{lib-micro}, an eBPF-inspired C macro language for
chaining $\mu$ops into triads together with a C API that installs a patch and
programs a match entry on an individual core, writing the sequencer arrays
through the \texttt{udbgrd}/\texttt{udbgwr} interface rather than through the
signed boot-time update path. Reaching that interface requires a part held in the
Red-Unlock state above, which we reproduce on our \hbox{Celeron~N3350}. Where
this body of work used such access to study or to harden the CPU, we use it to
make a cryptographic kernel the body of a hooked instruction, and the rest of the
paper is concerned with doing so well. The two ways in which a primitive
stresses the mechanism, control flow on the one hand and arithmetic density on
the other, are the subjects of \Cref{sec:keccak} and \Cref{sec:curve25519}
respectively.

% ---------------------------------------------------------------------
\subsection{The Triad as an Issue Bundle and the Entry Path}
\label{sec:triads-packing}

\paragraph{The entry path.}
Every primitive shares a single entry path that instantiates the customisation
mechanism of \Cref{sec:triad-anatomy}. We hook an otherwise unused
macroinstruction, the \texttt{vmwrite} whose opcode is \texttt{0x0cd8}, by
programming a match register so that its MSROM entry point is redirected into
our patch in the MSRAM. The native side places the operand pointers in fixed
architectural registers, with a base pointer in \texttt{RCX}, and issues a
single \texttt{vmwrite}, whereupon control transfers into the patch. The patch
reads its inputs, computes entirely in registers, writes its outputs through
the base pointer, and returns when its final triad carries the sequence word
that signals decode-complete. Because the hook is placed on a single static
instruction, a complete operation executes within one firing, which, as
\Cref{sec:budget} explains, is the property that makes register residency
worthwhile. When two patches must be resident simultaneously we hook a distinct
macroinstruction for each, since a single hooked instruction can be redirected
to only one entry point; the Curve25519 field operations of
\Cref{sec:curve25519}, for example, fire the multiplication on \texttt{vmwrite}
and the squaring on \texttt{vmread}, whose opcode is \texttt{0x0618}.

\paragraph{The triad as an issue bundle.}
Where \Cref{sec:triad-anatomy} described the triad as the microcode store sees
it, we now describe how we use it. We program each triad as a bundle of up to
three $\mu$ops together with its sequence word, and we rely throughout on a
property that we verified rather than assumed, as reported in
\Cref{sec:findings}: the three slots execute with full sequential semantics
over the register file. A value written in an earlier slot is visible to later
slots, which is read-after-write forwarding; a read in an earlier slot observes
the old value of a register that a later slot overwrites, which is the
write-after-read case; and when two slots write the same register the later
slot prevails, which is the write-after-write case. This holds identically for
general-purpose and temporary registers, and for condition flags as well as for
data. In practice the triad behaves like a small statically scheduled VLIW word
with intra-bundle forwarding, so that a short dependent chain can be folded into
a single issue group rather than spread across three instructions while
independent work fills the remaining slots. The layer is most expressive
precisely where compiled native code is weakest, namely in the dependent
reductions and carry chains that the instruction set architecture would
serialise; this is the lever behind both worked examples, where
\Cref{sec:keccak} folds a balanced parity tree into successive slots and
\Cref{sec:curve25519} fuses a multiplication, an accumulation and a carry
capture into one bundle.

\paragraph{Packing constraints.}
Two rules bound what may share a triad, and both are enforced by the hardware on
pain of a hang or a silently incorrect result. The first is that at most one
memory $\mu$op may occupy a triad, because a second load or store in the same
triad hangs the core, which has only a single first-level-cache data port. The
second is that a value produced by a load in an earlier slot is not forwarded to
a later slot of the same triad, unlike a register result, so that a consumer of
a freshly loaded value must begin a new triad. The asymmetry is worth stating
explicitly, since arithmetic-to-arithmetic forwarding within a triad is
permitted whereas load-to-arithmetic forwarding is not, and it is this boundary
condition that most often terminates a triad.

\paragraph{A self-verifying generator.}
Because a single mis-scheduled $\mu$op produces a silently incorrect result and
every hardware trial risks hanging the machine, we do not author triads by hand
where we can avoid it. For each primitive a generator emits an ordered op-list
over real register names, simulates that op-list against a model of the register
file and memory that incorporates the addressing constraints of
\Cref{sec:findings}, and asserts the result against a reference implementation
over many random inputs and against published known-answer vectors. Only
thereafter does a packer lower the op-list to triads and the generator emit the
C array. The packer is a single greedy left-to-right pass that closes the
current triad whenever it fills its three slots, whenever a second memory
operation would enter it, or whenever an operation would consume a register
loaded earlier in the same triad, so that it mechanically enforces the two
packing rules above, and it isolates control-transfer triads where required by
\Cref{sec:budget}. This discipline moves all correctness debugging into
software, where it is immediate and safe, and converts the hardware-specific
scheduling rules into a property of the toolchain rather than a burden on the
programmer. The same procedure of generation, simulation, verification and
packing underlies the Keccak permutation of \Cref{sec:keccak}; the small,
densely irregular field kernels of \Cref{sec:curve25519} are instead
hand-scheduled, but held to the identical verification discipline.

% ---------------------------------------------------------------------
\subsection{Patch RAM Budget}
\label{sec:budget}

The writable patch RAM on Goldmont holds 128 triads, occupying the region
\texttt{U7c00}--\texttt{U7dfc} in the MSRAM, in which each triad occupies four
address units and must begin on an even address, and within which a small
sub-range is reserved as a staging area and is not available to a primitive.
This ceiling of 128 triads is the binding resource of the entire approach, and
it shapes the design at the level of the algorithm rather than merely at the
level of the code.

\paragraph{The budget selects the strategy.}
Because a patch is small, the central question for each primitive is how to
accommodate a complete operation, including any iteration it involves, within
128 triads while keeping its working set resident. The answer is
primitive-specific but is in every case a budgetary decision. An operation whose
natural straight-line form would overflow the cap is restructured so that a
single resident body is reused: an iterated computation is looped in place by
means of a backward branch in the sequence word, with loop counters and constant
tables held in the operand buffer rather than encoded in the program, so that
the triad count becomes independent of the iteration count. This is the route
taken by Keccak in \Cref{sec:keccak}, whose 24 rounds are looped over one
resident round body. Where looping is too costly relative to the work performed
per iteration, we instead fold a fixed and small number of iterations into one
firing, as in the processing of two reduction words per \texttt{vmwrite} in the
word-by-word Montgomery primitives, and re-enter from the native side, thereby
trading a little additional input and output for remaining within the cap. By
contrast, a primitive whose straight-line form already fits, as the Curve25519
field operations of \Cref{sec:curve25519} do at 42 and 66 triads, is bound not
by the cap but by how densely its arithmetic can be packed. In every case it is
the relation between the work and the cap, rather than the instruction mix, that
fixes the macro-structure of the implementation.

\paragraph{Spend the budget on residency, not on movement.}
Temporary registers do not survive across firings, so any value that must
outlive a hooked instruction has to be carried in an architectural register or
spilled to memory. This makes a single firing per operation the default
arrangement, in which the inputs are loaded once in a prologue, the operation is
computed entirely in registers, and the outputs are stored once in an epilogue,
so that the memory traffic is paid once rather than once per iteration. The
doubled register file is what renders this feasible, and we regard every triad
spent on address arithmetic or data movement as a triad taken from the primitive
itself. A concrete illustration arises in our own work, in which a variant that
rebuilt an addressing base on each iteration in order to free one register
measured strictly slower than a variant that kept the base pinned and instead
borrowed a normally reserved register as scratch, because the cost of the rebuild
exceeded the cost of the spill it removed. Under a fixed triad budget, the
minimisation of data movement is the optimisation that matters most.

% ---------------------------------------------------------------------
\subsection{Findings}
\label{sec:findings}

Using microcode as a target required us first to establish the undocumented
behaviour of this particular Goldmont part, namely which $\mu$ops exist and what
they compute, how operands and flags are encoded, and which combinations are
illegal. Each finding reported below was isolated by a dedicated probe and is
now encoded as a rule in the generator and packer, so that a software-verified
op-list lowers to a patch that runs correctly on its first hardware trial. We
group the findings under the headings of the $\mu$op repertoire, the execution
and flag semantics, the addressing, the control flow, and the performance.

\paragraph{Micro-op repertoire.}
The vertically encoded $\mu$ops expose a RISC-like set that is sufficient for
field arithmetic, comprising 64-bit \texttt{ADD}, \texttt{XOR}, \texttt{SHL}
and \texttt{SHR}, the rotation \texttt{ROL}, which is a correct single $\mu$op
for every rotation amount, the bitwise operation
\texttt{NOTAND}$(d,a,b)=(\lnot a)\land b$, register moves by means of
\texttt{ZEROEXT}, and a $64\times64\rightarrow128$ multiplication
\texttt{MUL\_DSZ64\_DRR}$(hi,a,b)$ that preserves $a$, returns the low half in
$b$, and places the high half in $hi$. Immediate operand fields are 16 bits
wide, so a 64-bit constant must be materialised by composition, using
\texttt{ZEROEXT} followed by \texttt{SHL} and \texttt{XOR}, rather than in a
single operation. Two pitfalls cost us considerable time. The first is that the
memory operand forms, such as \texttt{ZEROEXT\_DSZ64\_DM}, do not perform
general addressing but instead read a field of the triggering macroinstruction,
so that general memory access is available only through the dedicated load and
store operations. The second is that operations with an implicit memory operand
cannot be used inside a patch at all. We therefore restrict the generator to
register-only arithmetic together with the explicit load and store operations.

\paragraph{Execution and flag semantics.}
Intra-triad register dataflow is fully sequential, as described in
\Cref{sec:triads-packing}, with the single exception that the result of a load
does not forward within a triad; store-to-load forwarding across triads within a
patch does function, however, which is what makes the spilling and reloading of
scratch values safe. A 64-bit multiplication may occupy any slot of a triad, an
earlier belief in a slot-zero restriction having not survived testing. The flags
are the subtlest aspect of the machine, and we found two distinct flag domains.
The first is a per-temporary-register condition that is set by an \texttt{ADD}
and consumed by a conditional \texttt{SETCC}, and that persists until the next
\texttt{ADD} to that register. The second is the architectural \texttt{RFLAGS},
which is frozen at patch entry and which is what \texttt{ADC} and the conditional
branch read. The bridge from the first domain to the second is unreliable in
large patches. Three practical consequences follow and are encoded in the
generator. All multi-word carry chains are built from per-register \texttt{ADD}
and \texttt{SETCC} pairs rather than from \texttt{ADC}; \texttt{SETCC} writes
correctly only to a temporary register and never to an architectural one, so a
carry must be materialised in a temporary and then moved out; and a temporary
needed across iterations must never be clobbered by an intervening carry
computation. The first of these is precisely what enables the parallel carry
chains of \Cref{sec:curve25519}, and its boundary is precisely what limits them
on the saturated representation discussed there.

\paragraph{Addressing.}
Loads and stores function only with a specific data-segment selector on this
part, and a selector copied from unrelated code hung the core unrecoverably, so
the working value is machine-specific and is fixed once. The immediate
displacement is an 8-bit signed byte offset spanning $\pm 128$ bytes, and an
early layout that placed scratch beyond that range wrapped silently and
corrupted state, a fault at first misdiagnosed as a dependency error. We
therefore centre the operand buffer beneath its base pointer so that all
displacements remain in range, and we use the register-indexed addressing form,
which carries no such limit, for the tables and counters that lie beyond the
reach of the immediate. As noted above, at most one memory operation may occupy
a triad.

\paragraph{Control flow.}
A backward conditional branch encoded in the sequence word transfers control to
an arbitrary triad and reads the architectural flags, ignoring any register
operand encoded alongside it, and the arithmetic operation that sets the tested
flag must occupy the same triad as the branch. This is the mechanism that
underlies the in-place looping of \Cref{sec:budget}, and \Cref{sec:keccak}
depends on it directly. Because temporary registers do not persist across
firings, any state that must survive between firings has to reside in an
architectural register or in memory, which is precisely why a complete, iterated
operation is kept within a single firing wherever the triad budget permits.

\paragraph{Performance.}
The layer has two regimes of per-triad cost, and conflating them is the easiest
way to misjudge a projected result. Measured as throughput, across many
independent firings that the out-of-order front end is able to overlap, a triad
costs approximately 0.49 cycles. A single dependent operation, by contrast, pays
the latency, and in the absence of any overlap between firings and in the
presence of a per-iteration branch the effective cost rises to roughly one triad
per cycle. A primitive is therefore fast in this layer to the extent that its
triad count is low and its critical path short, rather than because individual
operations are cheap. The improvements we report across the primitives, which
match or exceed both GCC at \texttt{-O3} and a state-of-the-art superoptimiser on
most operations, derive from this structural source: the doubled register file
removes the spills that dominate compiled code on an in-order core, and the
three-wide hazard-aware packing shortens dependent chains. The advantage is
structural, deriving from the elimination of data movement that the instruction
set architecture cannot avoid, rather than from a higher peak throughput; and
because the 128-triad cap forecloses extensive unrolling, the margin is modest on
kernels that are already tight and latency-bound and large on kernels that are
spill-heavy. The two worked examples that follow are chosen to exhibit both ends
of this range.
